\documentclass[12pt]{beamer}
\usetheme{Madrid}
\usecolortheme{whale}

\usepackage{graphicx}
\usepackage{booktabs}
\usepackage{amsmath}
\usepackage{amsfonts}
\usepackage{tikz}
\usetikzlibrary{calc}

% --- Presentation Details ---
\title{AI-Powered Job Application Automation Platform}
\author{
\footnotesize
\textbf{Presented by Group 10}\\
Aditya Sahu (CSE-22/01) \\
Aman Haque (CSE-22/02) \\
Ayush Kumar Ray (CSE-22/07) \\
Md. Noor Ali (CSE-22/22)
}

\institute{
Department of Computer Science and Engineering\\
DUIET, Dibrugarh University\\[0.5em]
Under the Supervision of\\
\textbf{Dr. Abhijit Boruah}\\
Assistant Professor, DUIET
}

\setbeamerfont{author}{size=\footnotesize}
\setbeamerfont{institute}{size=\small}

\setbeamertemplate{footline}{
  \leavevmode
  \hbox{
    \begin{beamercolorbox}[wd=.90\paperwidth,ht=2.5ex,dp=1ex,left]{title in head/foot}
      \hspace{0.3cm}\insertshorttitle
    \end{beamercolorbox}
    \begin{beamercolorbox}[wd=.10\paperwidth,ht=2.5ex,dp=1ex,right]{date in head/foot}
      \insertframenumber/\inserttotalframenumber\hspace{0.3cm}
    \end{beamercolorbox}
  }
}
\begin{document}

\begin{frame}[plain]

\centering

\begin{beamercolorbox}[
    wd=0.95\paperwidth,
    sep=8pt,
    center,
    rounded=true,
    shadow=true
]{title}
    \usebeamerfont{title}
    \color{white}
    \inserttitle
\end{beamercolorbox}

\vspace{0.25cm}

\includegraphics[width=2cm]{dibrugarh_logo.png}

\vspace{0.15cm}

{\usebeamerfont{author}\insertauthor\par}

\vspace{0.1cm}

{\usebeamerfont{institute}\insertinstitute\par}

\end{frame}


% --- Global Font Size Enforcement ---
\setbeamerfont{itemize/enumerate body}{size=\normalsize}
\setbeamerfont{itemize/enumerate subbody}{size=\normalsize}
\setbeamerfont{block title}{size=\normalsize}
\setbeamerfont{block body}{size=\normalsize}

\begin{document}

% --- Table of Contents Slide ---
\begin{frame}{Table of Contents}
    \tableofcontents
\end{frame}

\section{Introduction}

% --- Slide 1: Introduction & Problem Statement ---
\begin{frame}[allowframebreaks]{Introduction \& Problem Statement}
    \begin{block}{Introduction}
        \begin{itemize}
            \item Modern job seeking demands complex multi-step workflows: generating tailored resumes, crafting personalized cover letters, and setting up cold outreach.
            \item Pre-trained LLMs lack optimization for domain-specific professional formatting.
            
        \end{itemize}
    \end{block}
    
    \framebreak 
    
    \begin{block}{The Problem Statement}
        \begin{itemize}
            \item \textbf{System Constraints:} Existing resume builders are rigid, while traditional AI tools rely on expensive commercial APIs or fragile prompt engineering, resulting in unpredictable layouts.
            \item \textbf{Job Search Inefficiency:} Candidates must manually filter roles across multiple portals based on salary, tech stack, and qualifications. Scattered formats make comparisons difficult.
            \item \textbf{High Overhead:} Tailoring ATS-friendly resumes for specific job descriptions is heavily repetitive, cognitively taxing, and prone to human error.
            \item \textbf{Outreach Velocity:} Manually gathering profile links significantly slows down application momentum.
        \end{itemize}
    \end{block}
\end{frame}

\section{Literature Review}

\begin{frame}{Literature Review (1/2)}
\footnotesize

\begin{itemize}

\item \textbf{Resume Screening Using LLMs} [1] introduced semantic candidate ranking using vector embeddings and similarity matching with job descriptions.

\item \textbf{AI-Powered Resume Evaluation System} [2] focused on ATS compatibility, resume quality assessment, and improvement recommendations.

\item \textbf{LLM Agents in Recruitment} [3] proposed agent-based resume classification, candidate grading, and automated recruitment support.

\item \textbf{ResumeFlow} [4] utilized GPT-4 and Gemini for personalized resume refinement and job-specific customization.

\item \textbf{Professional Resume Generator} [5] employed BERT-based techniques for ATS optimization and resume enhancement.

\end{itemize}

\vspace{0.2cm}



\end{frame}

\begin{frame}{Literature Review (2/2)}
\footnotesize

\begin{itemize}

\item \textbf{CareerCorpus} [6] and \textbf{ResumeBench} [8] provided datasets and benchmarks for resume parsing and information extraction.

\item \textbf{Synthetic Resume Generation} [7] demonstrated the effectiveness of LLM-generated resumes for dataset augmentation.

\item \textbf{Mistral-7B} [9] introduced an efficient transformer architecture with strong language generation capabilities.

\item \textbf{QLoRA} [10] enabled memory-efficient fine-tuning through 4-bit NF4 quantization and low-rank adaptation.

\item \textbf{LangGraph} [11] provided a framework for building stateful multi-agent workflows suitable for automation pipelines.

\end{itemize}

\vspace{0.2cm}



\end{frame}

\section{Research Gap}

% --- Slide 2: Research Gap ---
\begin{frame}{Research Gap}
    \begin{block}{Research Gap}
        \begin{itemize}
            \item Existing literature primarily targets the employer-side (automated resume screening, parsing, and candidate ranking frameworks).
            \item There is a critical lack of open-source, fine-tuned localized language models dedicated to generating structured, compilable LaTeX application documents.
            \item No unified infrastructure exists that combines custom-trained models with multi-agent pipelines to enable autonomous job discovery, resume tailoring, and one-click job application.
        \end{itemize}
    \end{block}
\end{frame}

\section{Objectives}

% --- Slide 2b: Project Objectives ---
\begin{frame}{Project Objectives}

\small

\begin{block}{Project Objectives}
\begin{itemize}
    \item Automate job application.
    \item Generate ATS-friendly resumes.
    \item Create personalized cover letters and outreach emails.
    \item Recommend suitable jobs.
    \item Provide a complete one-click application platform using a fine-tuned Mistral-7B model.

\end{itemize}
\end{block}

\end{frame}

\section{Dataset}

\begin{frame}{Dataset Overview}
\small
\begin{itemize}
    \item Custom instruction-tuning dataset for ATS-friendly resume generation.
    
    \item Contains \textbf{10,357} training samples in JSONL format.
    
    \item Hybrid dataset:
    \begin{itemize}
        \item \textbf{40\%} real-world resumes
        \item \textbf{60\%} synthetic resumes
    \end{itemize}
    
    \item Each sample contains:
    \begin{itemize}
        \item \textbf{Instruction}: Resume generation task
        \item \textbf{Input}: Candidate profile information
        \item \textbf{Output}: ATS-friendly LaTeX resume
    \end{itemize}
    
    \item Domains covered:
    \begin{itemize}
        \item Software Engineering, Cloud Computing \& DevOps
        \item Web Development
        \item Data Science
        \item Machine Learning
    \end{itemize}
\end{itemize}
\end{frame}
\begin{frame}{Dataset Preparation Process}
\small

\begin{columns}[T]

% Left Side
\column{0.67\textwidth}

\begin{enumerate}
    \item Generated and collected structured candidate profiles.

    \item Standardized resume sections such as Education, Experience, Projects, Skills, and Achievements.

    \item Converted candidate information into instruction-response training pairs.

    \item Generated ATS-friendly resumes using a predefined LaTeX template.

    \item Removed incomplete samples and validated formatting consistency.

    \item Stored the final dataset in JSONL format for supervised fine-tuning.
\end{enumerate}

% Right Side
\column{0.33\textwidth}

\footnotesize
\textbf{Training Record Structure}

\vspace{0.2cm}

\ttfamily
\scriptsize

\{
\vspace{0.05cm}

\hspace*{0.2cm}"instruction": "Generate ATS-friendly resume",

\vspace{0.05cm}

\hspace*{0.2cm}"input": "Candidate profile details",

\vspace{0.05cm}

\hspace*{0.2cm}"output": "LaTeX resume code"

\vspace{0.05cm}

\}

\end{columns}

\end{frame}

\section{Model Architecture}

\begin{frame}[plain]
\frametitle{Mistral 7B Architecture and Implementation}

\centering

\includegraphics[
width=\paperwidth,
height=0.9\paperheight,
keepaspectratio
]{model_architecture.jpeg}

\end{frame}



\begin{frame}{Mistral-7B Fine-Tuning}
\small

\begin{block}{Tokenization using Mistral Tokenizer}
\footnotesize
\begin{itemize}
    \item Converted resume text into numerical token IDs using the Mistral tokenizer.
    
    \item Applied subword tokenization to efficiently handle technical terms, skills, and domain-specific vocabulary.
    
    \item Generated input token sequence:
    \[
    X = [x_1, x_2, x_3, \ldots, x_n]
    \]
    
    \item Created attention masks to distinguish valid tokens from padding tokens.
    
    \item Produced input-output tensor pairs for supervised fine-tuning.
    
    \item Enabled efficient processing of long resume documents while preserving contextual information.
\end{itemize}
\end{block}

\vspace{0.2cm}

\textbf{Tokenization Workflow}

\[
\text{Resume Text}
\rightarrow
\text{Tokenizer}
\rightarrow
\text{Token IDs}
\rightarrow
\text{Tensors}
\rightarrow
\text{Mistral-7B}
\]

\end{frame}

\section{Fine-Tuning Methodology}

% --- Slide 4: ML Training & Implementation ---
\begin{frame}{QLoRA Quantization Strategy}
\vspace{-0.2cm}
\footnotesize
\begin{columns}[T]

% Left Column
\column{0.45\textwidth}
\begin{itemize}\setlength\itemsep{0.05em}
    \item Applied \textbf{QLoRA} to fine-tune Mistral-7B efficiently.
    \item Quantized model weights to \textbf{4-bit NF4}, reducing VRAM requirements.
    \item Kept pretrained weights frozen and trained only low-rank adapter matrices.
    \item Reduced memory usage from \textbf{112 GB} to \textbf{14 GB}.
    \item Achieved \textbf{87.5\% VRAM reduction} while maintaining model performance.
\end{itemize}
\vspace{0.08cm}
\[
\text{4-bit NF4} \rightarrow \text{LoRA} \rightarrow \text{Fine-Tuning}
\]
\vspace{0.05cm}
{\scriptsize\textbf{Performance Comparison}}\\[0.1cm]
{\tiny
\begin{tabular}{@{}lcc@{}}
\hline
\textbf{Metric} & \textbf{Full FT} & \textbf{QLoRA} \\
\hline
VRAM Usage     & 112 GB   & 14 GB   \\
ROUGE-L        & 0.61     & 0.59    \\
Training Time  & $\sim$8h & $\sim$3h \\
Output Quality & High     & Comparable \\
\hline
\end{tabular}
}

% Right Column
\column{0.55\textwidth}
\centering
\includegraphics[width=\linewidth,height=0.38\textheight,keepaspectratio]{qlora_vram.png}
\\[0.05cm]
{\scriptsize VRAM comparison between Full Fine-Tuning and QLoRA.}
\\[0.1cm]
\textbf{LoRA Update Rule}
\[
W' = W + BA
\]
{\scriptsize
$W$ = Frozen pretrained weights \qquad
$A,B$ = Trainable low-rank matrices
}
\end{columns}
\end{frame}

\begin{frame}{LoRA Adaptation}
Low-rank trainable matrices injected into linear modules instead of full-parameter updates:
\[ W' = W + \Delta W = W_q + \frac{\alpha}{r}(B \cdot A) \]
\begin{itemize}
  \item $W \in \mathbb{R}^{d\times k}$: frozen base weight
  \item $B \in \mathbb{R}^{d\times r}$, $A \in \mathbb{R}^{r \times k}$: trainable low-rank matrices ($r \ll \min(d,k)$)
  \item $\alpha$: scaling hyperparameter
\end{itemize}

\begin{table}
\centering
\small
\begin{tabular}{ll}
\toprule
\textbf{Parameter} & \textbf{Value} \\
\midrule
Epochs & 10 \\
Batch Size / Grad. Accum. & 2 / 8 \\
Learning Rate & $2\times10^{-5}$ \\
LoRA Rank $r$ / Alpha $\alpha$ & 16 / 32 \\
LoRA Dropout & 0.05 \\
Optimizer & Paged AdamW (32-bit) \\
\bottomrule
\end{tabular}
\end{table}
\end{frame}

\section{Results and Evaluation}

\begin{frame}{Results and Evaluation}
\footnotesize
\begin{columns}[T]

% Left Column - Metrics Explanation
\column{0.58\textwidth}
\textbf{Evaluation Metrics}\\[0.15cm]
\begin{itemize}\setlength\itemsep{0.4em}
    \item \textbf{BLEU}\\[0.05cm]
    {\scriptsize Evaluates generated text by comparing it with human-written reference text based on n-gram overlap.}

    \item \textbf{ROUGE-L}\\[0.05cm]
    {\scriptsize Measures the quality of generated text using the longest common subsequence between generated and reference content.}

    \item \textbf{BERTScore}\\[0.05cm]
    {\scriptsize Measures semantic similarity between generated and reference text using contextual embeddings from transformer models.}
\end{itemize}

% Right Column - Table
\column{0.42\textwidth}
\centering
\textbf{Resume Generation Performance}\\[0.2cm]
\begin{tabular}{|l|c|}
\hline
\textbf{Metric} & \textbf{Score} \\
\hline
BLEU          & 0.3921 \\
ROUGE-1       & 0.6908 \\
ROUGE-2       & 0.6859 \\
ROUGE-L       & 0.6904 \\
BERTScore F1  & \textbf{0.9542} \\
\hline
\end{tabular}

\vspace{0.25cm}
\raggedright
\textbf{}\\[0.1cm]
\begin{itemize}\setlength\itemsep{0.05em}
    \item BERTScore F1 of \textbf{0.9542} indicates high semantic similarity between generated and reference resumes.
\end{itemize}

\end{columns}
\end{frame}

\begin{frame}{Comparative Analysis with Existing Research}

\centering

\includegraphics[
width=0.98\textwidth,
height=1.0\textheight,
keepaspectratio
]{comparative_analysis.png}

\end{frame}

% ============================================================
% NEW SECTION: System Implementation
% ============================================================

\section{System Implementation}

\begin{frame}{Platform Architecture \& Tech Stack (1/2)}
\small
\begin{block}{Frontend Architecture}
\begin{itemize}\setlength\itemsep{0.3em}
    \item \textbf{Next.js 16 (React 19):} Server-Side Rendering (SSR) for fast initial load times and optimized routing.
    \item \textbf{TypeScript:} Enforces strict type safety across components, reducing runtime errors.
    \item \textbf{Tailwind CSS \& UI Components:} Delivers a highly responsive, modern, and accessible user interface using utility-first styling and custom animated components.
    \item \textbf{OAuth 2.0 Integration:} Secure, frictionless user authentication via Google and GitHub.
\end{itemize}
\end{block}

\vspace{0.2cm}

\begin{block}{Backend Architecture}
\begin{itemize}\setlength\itemsep{0.3em}
    \item \textbf{Node.js + Express.js:} High-performance RESTful API to handle client requests and route logic.
    \item \textbf{Prisma ORM + PostgreSQL:} Robust data modeling and structured persistence layer for user profiles, resumes, and job history.
    \item \textbf{JWT Authentication:} Stateless, scalable, and secure session management.
    \item \textbf{Inngest:} Serverless event-driven scheduling for reliable background job execution.
\end{itemize}
\end{block}
\end{frame}

\begin{frame}{Platform Architecture \& Tech Stack (2/2)}
\small
\begin{block}{AI \& Orchestration Layer}
\begin{itemize}\setlength\itemsep{0.3em}
    \item \textbf{LangGraph:} Manages complex, stateful multi-agent workflows, orchestrating the step-by-step logic of job application generation.
    \item \textbf{LangChain:} Abstracts LLM interactions and creates dynamic prompt chains for data extraction and document drafting.
    \item \textbf{Gemini Flash / Pro:} Acts as the primary inference engine for parsing job descriptions, ranking roles, and drafting cover letters/emails.
    \item \textbf{Fine-tuned Mistral-7B:} Specialized model exclusively deployed for structured, ATS-compliant LaTeX resume generation.
\end{itemize}
\end{block}

\vspace{0.2cm}

\begin{block}{Job Discovery \& Aggregation}
\begin{itemize}\setlength\itemsep{0.3em}
    \item \textbf{JSearch API (Primary):} Aggregates real-time job listings from major platforms like Google Jobs, LinkedIn, and Indeed.
    \item \textbf{Remotive API (Fallback):} Provides supplementary listings focused specifically on remote and distributed roles.
\end{itemize}
\end{block}
\end{frame}

\begin{frame}{Multi-Agent Pipeline Design: One-Click Apply}
\small

\begin{block}{Pipeline Workflow (LangGraph)}
\centering
\vspace{0.15cm}
\footnotesize
\texttt{parseJob} $\rightarrow$ \texttt{scoreMatch} $\rightarrow$ \texttt{tailorResume} $\rightarrow$ \texttt{generateCoverLetter} $\rightarrow$ \texttt{draftColdEmail}
\vspace{0.1cm}
\end{block}

\vspace{0.2cm}

\textbf{Node Functions \& Asset Generation}
\begin{itemize}\setlength\itemsep{0.4em}
    \item \textbf{parseJob:} Extracts critical skills, requirements, and keywords from the raw job description using structured LLM output.
    \item \textbf{scoreMatch:} Computes a quantitative ATS match score (0--100) by comparing the candidate's master profile against parsed job requirements.
    \item \textbf{tailorResume:} Synthesizes the parsed requirements and master profile to generate a highly targeted, compilable \LaTeX{} resume.
    \item \textbf{generateCoverLetter:} Drafts a personalized 300-word cover letter addressing specific pain points mentioned in the job post.
    \item \textbf{draftColdEmail:} Composes a concise, punchy outreach email optimized for recruiter engagement.
\end{itemize}

\end{frame}

\begin{frame}{Multi-Agent Pipeline Design: Job Autopilot}
\small

\begin{block}{Pipeline Workflow (Scheduled Execution via Inngest)}
\centering
\vspace{0.15cm}
\footnotesize
\texttt{buildSearchQuery} $\rightarrow$ \texttt{searchJobs} $\rightarrow$ \texttt{filterAndRank} $\rightarrow$ \texttt{batchApply} $\rightarrow$ \texttt{generateReport}
\vspace{0.1cm}
\end{block}

\vspace{0.2cm}

\textbf{Node Functions \& Automation Logic}
\begin{itemize}\setlength\itemsep{0.4em}
    \item \textbf{buildSearchQuery:} LLM analyzes the user profile to formulate 2--3 highly focused, board-optimized search queries.
    \item \textbf{searchJobs:} Executes queries across JSearch \& Remotive APIs, deduplicates listings, and limits results to maintain efficiency.
    \item \textbf{filterAndRank:} LLM evaluates and scores each discovered job against the master profile, filtering out low-match roles (score $<$ 70).
    \item \textbf{batchApply:} Iteratively triggers the One-Click Apply pipeline for every qualified job discovered in the batch.
    \item \textbf{generateReport:} Compiles a structured, comprehensive summary report detailing all jobs analyzed and application assets generated.
\end{itemize}

\end{frame}

% ============================================================
% END NEW SECTION
% ============================================================

\section{Conclusion}

\begin{frame}{Conclusion}

\small

\begin{itemize}
    \item Successfully developed an end-to-end \textbf{AI-powered job automation platform} that seamlessly bridges the gap between job discovery and application asset generation.
    \item The \textbf{One-Click Apply} pipeline dramatically reduces the cognitive load on candidates by autonomously tailoring \LaTeX{} resumes, cover letters, and cold outreach emails for any given job description.
    \item The \textbf{Job Autopilot} subsystem introduces fully scheduled, background job discovery, utilizing LLMs to dynamically formulate search queries, rank relevance, and execute batch applications without human intervention.
    \item By integrating a robust \textbf{LangGraph multi-agent architecture} with a modern Next.js/Node.js stack, the system offers high scalability, reliability, and an intuitive user experience.
    \item While leveraging a highly efficient \textbf{fine-tuned Mistral-7B} for specialized resume generation (achieving a BERTScore F1 of 0.9542), the platform's true innovation lies in its orchestration of multiple APIs and agentic workflows.
    \item \textbf{Future work} includes incorporating user feedback loops, expanding to additional job board integrations, and exploring real-time tracking of application submission statuses.
\end{itemize}

\end{frame}

\begin{frame}[allowframebreaks,noframenumbering]{References}
\scriptsize

\begin{thebibliography}{99}

\bibitem{ref1}
E. Salakar, Y. Shah, J. Rai, J. Wadmare, and A. Salian,
``Resume Screening Using Large Language Models,''
ICAST, IEEE, 2023.

\bibitem{ref2}
G. Hendre, S. Bhiware, A. Darwajkar, and A. S. Mali,
``AI-Powered Resume Evaluation System,''
IJARSCT, 2025.

\bibitem{ref3}
C. Gan, Q. Zhang, and T. Mori,
``Application of LLM Agents in Recruitment: A Novel Framework for Resume Screening,''
arXiv, 2024.

\bibitem{ref4}
S. B. Zinjad et al.,
``ResumeFlow: An LLM-Facilitated Pipeline for Personalized Resume Generation and Refinement,''
SIGIR, 2024.

\bibitem{ref5}
M. V. Reddy et al.,
``Professional Resume Generator,''
Parul University, 2025.

\bibitem{ref6}
M. S. Chowdhury, A. F. Chowdhury, A. Banu, and R. Hossain,
``CareerCorpus: A Comprehensive Dataset of Annotated Resumes,''
Data in Brief, 2026.

\bibitem{ref7}
P. Skondras, P. Zervas, and G. Tzimas,
``Generating Synthetic Resume Data with Large Language Models for Enhanced Job Description Classification,''
Future Internet, 2023.

\bibitem{ref8}
Z. Ling, H. Zhang, J. Cui, Z. Wu, X. Sun, G. Li, and X. He,
``Beyond Human Labels: A Multi-Linguistic Auto-Generated Benchmark for Evaluating Large Language Models on Resume Parsing,''
EMNLP, 2025.

\bibitem{ref9}
A. Q. Jiang et al.,
``Mistral 7B,''
arXiv:2310.06825, 2023.

\bibitem{ref10}
T. Dettmers, A. Pagnoni, A. Holtzman, and L. Zettlemoyer,
``QLoRA: Efficient Finetuning of Quantized LLMs,''
2023.

\bibitem{ref11}
LangChain Inc.,
``LangGraph: Building Stateful, Multi-Actor Applications with LLMs,''
2024.

\bibitem{ref12}
Inngest Inc.,
``Inngest: Serverless Background Functions and Workflow Engine,''
2024.

\bibitem{ref13}
H. Touvron et al.,
``LLaMA: Open and Efficient Foundation Language Models,''
2023.

\bibitem{ref14}
P. Dubey, S. Yadav, and S. Gawande,
``Resume Generator Using AI,''
2023.

\bibitem{ref15}
H. Saiyad, J. Khan, A. A. Khan, and A. Pathan,
``AI-Powered Resume Builder: Automating Professional Resume Generation Using NLP and Machine Learning,''
IJMRSET, 2025.

\bibitem{ref16}
A. Dubey, N. Dhanda, A. Lovewanshi, and V. Yadav,
``Resume2Role: Intelligent Resume Parsing and Domain Prediction Using NLP and Classification Models,''
ITM Web of Conferences, 2026.

\bibitem{ref17}
S. T. Gopalakrishna and V. Varadharajan,
``Automated Tool for Resume Classification Using Semantic Analysis,''
IJAIA, 2019.

\bibitem{ref18}
R. Pal, S. Shaikh, S. Satpute, and S. Bhagwat,
``Resume Classification Using Various Machine Learning Algorithms,''
ITM Web of Conferences, 2022.

\bibitem{ref19}
A. Heakl, Y. Mohamed, N. Mohamed, A. Elsharkawy, and A. Zaky,
``R\'esum\'eAtlas: Revisiting Resume Classification with Large-Scale Datasets and Large Language Models,''
Procedia Computer Science, 2024.

\bibitem{ref20}
C. G. Onukwugha, C. I. Ofoegbu, O. B. Aliche, and C. U. Betrand,
``Resume Optimization Model Using Machine Learning Techniques,''
International Journal of Intelligent Information Systems, 2024.

\bibitem{ref21}
S. Efeoglu and A. Paschke,
``Fine-Tuning Large Language Models for Relation Extraction within a Retrieval-Augmented Generation Framework,''
XLLM, 2025.

\bibitem{ref22}
S. M. Wong, H. Leung, and K. Y. Wong,
``Efficiency in Language Understanding and Generation: An Evaluation of Four Open-Source Large Language Models,''
Research Square, 2024.

\bibitem{ref23}
A. Mullick, S. Bose, R. Saha, A. K. Bhowmick, A. Vempaty,
P. Goyal, N. Ganguly, P. Dey, and R. Kokku,
``Leveraging the Power of LLMs: A Fine-Tuning Approach for High-Quality Aspect-Based Summarization,''
arXiv, 2024.

\end{thebibliography}

\end{frame}

\begin{frame}{}
\centering
\vfill

{\Huge \textbf{Thank You}}

\vfill
\end{frame}
\end{document}
